\section{Extended results and protocol}
\label{app:extended}

\paragraph{All reconstruction contrasts.} Table~\ref{tab:reconcontrasts} reports every
pairwise loss contrast in the Q16 reconstruction study, across the breakout and rebound
families, the May--July and August--November cohort groups, and both the equal-day and pooled
weightings, in raw and standardized form. The pattern that the main text summarizes is
visible throughout the table. The augmented-minus-coarse (P$-$C) and rich-minus-coarse
(R$-$C) contrasts are small and frequently positive, so augmentation and rich inputs lose.
The predictor-minus-frequency contrasts show that no learned arm reliably beats the
book-free baseline. The frequency-minus-uniform contrast confirms that the
training-frequency baseline is itself a strong reference. The rich-minus-augmented (R$-$P)
contrast is small in every cell. This is why reconstruction quality does not translate into
a decision: the observed-rich predictor and its reconstruction do not separate at the
decision even where the reconstruction is accurate.

\begin{longtable}{llllrrr}
\caption{Q16 reconstruction study: all pairwise loss contrasts by family, cohort group, and
weighting (raw and standardized). Positive means the first predictor loses more. Source:
\texttt{experiments/results/e264\_group\_metrics.csv}.}
\label{tab:reconcontrasts}\\
\toprule
family & group & weight & contrast & $n$ & raw & standardized \\
\midrule
\endfirsthead
\toprule
family & group & weight & contrast & $n$ & raw & standardized \\
\midrule
\endhead
\bottomrule
\endfoot
\end{longtable}


\paragraph{Q17 cost/latency sensitivity.} Table~\ref{tab:q17ranges} reports, for each asset
and policy, the utility range swept across fixed later-date cost and latency scenarios,
together with the number of cells whose sign opposes the point estimate. The argmax policy is
consistently negative in absolute terms, the confidence policy collapses to abstention, and
the scheduling policy straddles zero with many opposing cells. The negative therefore does not
depend on a single chosen cost point. It is stable across the scenario sweep, which is the
reason we report ranges rather than a tuned scenario.

\begin{table}[h]\centering
\caption{Q17 later-date sensitivity: utility range across fixed cost/latency scenarios per
asset and policy, with the count of cells whose sign opposes the point estimate. Source:
\texttt{experiments/results/q17\_figure\_data.json}.}
\label{tab:q17ranges}
\begin{tabular}{llrrrrr}
\toprule
asset & policy & min & max & opposite & cells & macro-F1 \\
\midrule
BTC & argmax & 0.0017 & 4.3955 & 0 & 54 & 0.1799 \\
BTC & fixed\_confidence & -1.7188 & 0.0577 & 45 & 54 & 0.1799 \\
ETH & argmax & 0.0279 & 5.2932 & 0 & 54 & 0.1585 \\
ETH & fixed\_confidence & -0.7974 & 0.0039 & 48 & 54 & 0.1585 \\
\bottomrule
\end{tabular}
\end{table}


\paragraph{Mixture-versus-within decomposition.} Table~\ref{tab:decomp} decomposes each
reconstruction loss gap into a class-mixture term (driven by differences in realized class
frequencies across cohorts) and a within-class term (driven by per-class loss differences).
The mixture term carries most of any apparent movement, which is exactly why an aggregate over
heterogeneous cohorts can manufacture a depth ``advantage'' that vanishes under a fixed
weighting: the aggregate is confounded by the changing class mix, not by information. The
within-class terms, which isolate genuine predictive change, are small and do not favor the
augmented or rich predictors.

\begin{longtable}{llllrrr}
\caption{Q16 loss-gap decomposition into a class-mixture term and a within-class term, by
family, contrast, and metric. The mixture term localizes the earlier apparent depth
advantage. Source: \texttt{experiments/results/e264\_decompositions.csv}.}
\label{tab:decomp}\\
\toprule
family & contrast & metric & status & $\Delta$ & mix & within \\
\midrule
\endfirsthead
\toprule
family & contrast & metric & status & $\Delta$ & mix & within \\
\midrule
\endhead
\bottomrule
\endfoot
breakout & equal\_day & C & evaluable & 0.006682 & 0.020763 & -0.014081 \\
breakout & equal\_day & P & evaluable & 0.013920 & 0.021523 & -0.007603 \\
breakout & equal\_day & R & evaluable & 0.024499 & 0.022805 & 0.001694 \\
breakout & equal\_day & onehot-F & evaluable & -0.021897 & -0.021897 & 0.000000 \\
breakout & equal\_day & onehot-A & evaluable & -0.022745 & -0.022745 & 0.000000 \\
breakout & equal\_day & onehot-N & evaluable & 0.044642 & 0.044642 & 0.000000 \\
breakout & equal\_day & uniform & evaluable & -0.000000 & 0.000000 & -0.000000 \\
breakout & equal\_day & train-frequency & evaluable & 0.013284 & 0.013284 & 0.000000 \\
breakout & equal\_day & P-C & evaluable & 0.007238 & 0.000760 & 0.006478 \\
breakout & equal\_day & R-C & evaluable & 0.017817 & 0.002042 & 0.015775 \\
breakout & equal\_day & R-P & evaluable & 0.010579 & 0.001282 & 0.009297 \\
breakout & equal\_day & C-frequency & evaluable & -0.006602 & 0.007479 & -0.014081 \\
breakout & equal\_day & P-frequency & evaluable & 0.000636 & 0.008239 & -0.007603 \\
breakout & equal\_day & R-frequency & evaluable & 0.011215 & 0.009521 & 0.001694 \\
breakout & equal\_day & frequency-uniform & evaluable & 0.013284 & 0.013284 & 0.000000 \\
breakout & pooled & C & evaluable & 0.023057 & 0.035486 & -0.012429 \\
breakout & pooled & P & evaluable & 0.029897 & 0.036989 & -0.007092 \\
breakout & pooled & R & evaluable & 0.042561 & 0.039458 & 0.003103 \\
breakout & pooled & onehot-F & evaluable & -0.053973 & -0.053973 & 0.000000 \\
breakout & pooled & onehot-A & evaluable & -0.030975 & -0.030975 & 0.000000 \\
breakout & pooled & onehot-N & evaluable & 0.084948 & 0.084948 & 0.000000 \\
breakout & pooled & uniform & evaluable & -0.000000 & 0.000000 & 0.000000 \\
breakout & pooled & train-frequency & evaluable & 0.022671 & 0.022671 & -0.000000 \\
breakout & pooled & P-C & evaluable & 0.006839 & 0.001503 & 0.005337 \\
breakout & pooled & R-C & evaluable & 0.019504 & 0.003972 & 0.015532 \\
breakout & pooled & R-P & evaluable & 0.012664 & 0.002469 & 0.010195 \\
breakout & pooled & C-frequency & evaluable & 0.000386 & 0.012815 & -0.012429 \\
breakout & pooled & P-frequency & evaluable & 0.007226 & 0.014318 & -0.007092 \\
breakout & pooled & R-frequency & evaluable & 0.019890 & 0.016787 & 0.003103 \\
breakout & pooled & frequency-uniform & evaluable & 0.022671 & 0.022671 & 0.000000 \\
rebound & equal\_day & C & evaluable & 0.008829 & 0.008453 & 0.000375 \\
rebound & equal\_day & P & evaluable & 0.008006 & 0.009437 & -0.001431 \\
rebound & equal\_day & R & evaluable & 0.018077 & 0.014875 & 0.003202 \\
rebound & equal\_day & onehot-F & evaluable & -0.016623 & -0.016623 & 0.000000 \\
rebound & equal\_day & onehot-A & evaluable & -0.046166 & -0.046166 & 0.000000 \\
rebound & equal\_day & onehot-N & evaluable & 0.062788 & 0.062788 & 0.000000 \\
rebound & equal\_day & uniform & evaluable & 0.000000 & -0.000000 & -0.000000 \\
rebound & equal\_day & train-frequency & evaluable & 0.006565 & 0.006565 & -0.000000 \\
rebound & equal\_day & P-C & evaluable & -0.000823 & 0.000984 & -0.001806 \\
rebound & equal\_day & R-C & evaluable & 0.009249 & 0.006422 & 0.002827 \\
rebound & equal\_day & R-P & evaluable & 0.010071 & 0.005438 & 0.004633 \\
rebound & equal\_day & C-frequency & evaluable & 0.002264 & 0.001888 & 0.000375 \\
rebound & equal\_day & P-frequency & evaluable & 0.001441 & 0.002872 & -0.001431 \\
rebound & equal\_day & R-frequency & evaluable & 0.011512 & 0.008310 & 0.003202 \\
rebound & equal\_day & frequency-uniform & evaluable & 0.006565 & 0.006565 & -0.000000 \\
rebound & pooled & C & evaluable & 0.014667 & 0.008656 & 0.006011 \\
rebound & pooled & P & evaluable & 0.016578 & 0.010108 & 0.006470 \\
rebound & pooled & R & evaluable & 0.021834 & 0.017754 & 0.004080 \\
rebound & pooled & onehot-F & evaluable & -0.019597 & -0.019597 & 0.000000 \\
rebound & pooled & onehot-A & evaluable & -0.075886 & -0.075886 & 0.000000 \\
rebound & pooled & onehot-N & evaluable & 0.095482 & 0.095482 & 0.000000 \\
rebound & pooled & uniform & evaluable & -0.000000 & 0.000000 & 0.000000 \\
rebound & pooled & train-frequency & evaluable & 0.007775 & 0.007775 & -0.000000 \\
rebound & pooled & P-C & evaluable & 0.001911 & 0.001452 & 0.000458 \\
rebound & pooled & R-C & evaluable & 0.007166 & 0.009098 & -0.001931 \\
rebound & pooled & R-P & evaluable & 0.005256 & 0.007646 & -0.002390 \\
rebound & pooled & C-frequency & evaluable & 0.006892 & 0.000881 & 0.006011 \\
rebound & pooled & P-frequency & evaluable & 0.008803 & 0.002333 & 0.006470 \\
rebound & pooled & R-frequency & evaluable & 0.014058 & 0.009979 & 0.004080 \\
rebound & pooled & frequency-uniform & evaluable & 0.007775 & 0.007775 & -0.000000 \\
\end{longtable}


\paragraph{Protocol in pseudocode.} The matched intervention is deliberately simple, and its
value is in what it forbids. For each study we (1) freeze the model family, training budget,
and calendar split before any test date is read. (2) construct features and normalization on
the training portion only. (3) enumerate the information views to compare, changing nothing
else between them. (4) fit each view under the shared budget, preserving convergence warnings;
(5) evaluate every view on the identical held-out events, paired at the event level. (6) score
against a fixed metric and, for decisions, a fixed policy and explicit cost model. And (7)
apply the pre-registered gate without any outcome-driven adjustment. A failure to certify an
effect is recorded as such and never reported as evidence of the opposite effect. This is the
whole instrument, and the negatives in this paper are what it returns when it is not allowed
to shop for an outcome.
